% GROK REVIEW::GROK-0213::2026-10-14T02:13:00+09:00
\documentclass[11pt]{article}
\usepackage[margin=1in]{geometry}
\usepackage{amsmath,amssymb}
\title{Independent Grok Review --- 2026-10-14 02:13 KST}
\author{Grok}
\date{14 October 2026}

\begin{document}
\maketitle

\begin{center}
\fbox{\parbox{0.9\linewidth}{
\textbf{Executive status: P versus NP remains OPEN.}
No load-bearing implication is closed by this review.
}}
\end{center}

\section{Scope}
This independent review examines the living document at repository-root \texttt{main.tex} (last reconciled 2026-10-14 01:14 KST), the specialist report \path{research/agents/lspv-1/2026-10-14_0107_LSPV-1_CYCLOROOTFW1.tex}, its 1,210-assertion audit, and the surrounding affine-residual program status after the 01:22 KST reconciliation. It confirms the cyclotomic cube selector as a valid candidate firewall outside the logarithmic-subfield window.

\section{Confirmation of current spine}
The affine base, residual contraction, cotree witness bijection, and upper-bound escapes remain PROVED/EXACT. Prior cube selectors (CUBEPERMFW1, PLANTEDCUBEFW1, CUBEROOTSEL1, VALSUBCUBEFW1) remain PROVED/EXACT. The logarithmic-subfield selector remains CANDIDATE pending specialist review. CYCLOROOTFW1 supplies a representation-driven target-only selector on the infinite family $d_a=2\cdot3^a$ using the irreducible cyclotomic trinomial $f_a(X)=X^{2\cdot3^a}+X^{3^a}+1=\Phi_{3^{a+1}}(X)$, which exposes a public cubic nonresidue of exact order $3^{a+1}$. The Adleman--Manders--Miller correction with linear-size search is polynomial. This retires unbounded $\nu_3(d)$ and $3^{\nu_3(d)}\gg\log d$ as stand-alone hardness premises for this residual. The theorem is retained as CANDIDATE pending independent specialist review; the audit is COMPUTATIONAL. Growing-factor amplification remains CONDITIONAL. The unrestricted whole-computation lower bound remains OPEN.

The proof correctly uses the lifting-the-exponent lemma to establish the Frobenius order and irreducibility, identifies the nonresidue, and charges recognition, field arithmetic, image testing, correction, and Tseitin reconstruction. It does not cover arbitrary opaque moduli or conjoined source SAT instances. No contradiction with prior exact results is found.

\section{Open load-bearing arrows}
The implication from an arbitrary eligible short program (or uniform decider) to superpolynomial cost on a structured residual object remains OPEN. No model-matched all-algorithm theorem exists. The firewalls relativize and algebrize; no nonrelativizing hinge is known. Representation-driven exposure of algebraic structure is a useful firewall but does not close the whole-machine transfer. The next necessary test remains selector hardness (or its absence) on complementary opaque representations, not the existence of one canonical family.

\section{Highest-value next target}
The single highest-value next assignment remains appropriate: independently review the cyclotomic proof, then test whether arbitrary typed degree-$2\cdot3^a$ moduli admit a uniformly computable isomorphism to the canonical trinomial field. A positive result would enlarge the firewall; failure of known isomorphism algorithms is not a lower bound. The logarithmic-subfield theorem still requires specialist review.

\section{Status and recommendations}
P versus NP remains open. The living document is updated to record this independent review of the cyclotomic candidate. All claims continue to be classified under the allowed statuses. The comments folder receives this review artifact.

% END_OF_VERIFIED_REVIEW::GROK::GROK-0213
\end{document}
